\documentclass[11pt]{article}

\usepackage[utf8]{inputenc}
\usepackage[T1]{fontenc}
\usepackage{amsmath}
\usepackage{booktabs}
\usepackage[margin=1.15in]{geometry}
\usepackage[hidelinks]{hyperref}
\usepackage{enumitem}
\emergencystretch=2em
\setlength{\parskip}{.55em}
\setlength{\parindent}{0pt}

\title{How Lexodle scores a guess, in plain terms}
\author{The Lexodle project}
\date{\today}

\begin{document}
\maketitle

\textbf{The maths here is mostly AI-generated and has not been peer reviewed.}

This is the same material as the paper \texttt{main.pdf}, with the notation taken out. You need
fractions, percentages and averages. For the logarithm, one calculator button (log base 2) is
enough, and every number in the examples can be checked with it.

\tableofcontents

\section{The game, and what is being measured}

A word is hidden. You guess a word of the same length, and each letter comes back green (right
place), yellow (in the word, wrong place) or grey (not in the word). You keep guessing until you
find the word or run out of tries.

Every guess is a question and the colours are the answer. Before a guess some words are still
possible. After it, only the words that would have given \emph{exactly those colours} are still
possible. A good guess leaves few words, and the answer tells you a lot.

Lexodle measures that in \textbf{bits}. A bit is what it takes to cut the possible words in half.
Eight words left to one is three bits; 1,024 to one is ten. Finding the word means gaining all the
bits there are to find: for $N$ equally likely words, $\log_2 N$ of them.

\section{Bits}\label{sec:bits}

A guess splits the possible words into groups, one for each set of colours it could get. If you
land in a group that held a share $q$ of the chances, you \textbf{gain} $-\log_2 q$ bits. A group
holding half the chances gains one bit; one holding a quarter gains two; one holding a thousandth
gains about ten.

\textbf{Example.} Three words, equally likely, and a guess that tells apart one of them from the
other two. Groups: one word ($q=1/3$, gain $1.58$ bits) and two words ($q=2/3$, gain $0.58$ bits).

The \textbf{expected gain}, written $H$, is each outcome's gain times its chance, added up:
$\tfrac13\times1.58+\tfrac23\times0.58=0.92$ bits. It is what the guess is worth before you know
the colours.

\section{Skill}\label{sec:skill}

Every word the game accepts could have been your guess. The best of them has the largest expected
gain, $H^*$. Your \textbf{skill} is your guess's $H$ as a share of $H^*$, as a percentage.

In the example, suppose the best guess could gain $1.58$ bits (every word alone). Then the guess
above has skill $0.92/1.58\times100=58$. A guess that tells you nothing (one group) has skill 0.

Skill does not depend on the hidden word. It depends only on what was known and on the guess.

\section{Luck}\label{sec:luck}

The colours you got gained some number of bits. Luck says where that gain falls among the gains the
guess could have given: the chance of a smaller gain, minus the chance of a larger one, as a
percentage.

In the example, landing in the single word (gain $1.58$, the biggest, a $1/3$ chance) has luck
$2/3-0=+67$. Landing in the pair (gain $0.58$, a $2/3$ chance) has luck $0-1/3=-33$. Across all the
ways the answer could be, luck averages exactly 0: it is neither a reward nor a penalty, only how
the colours fell.

Luck is also shown in bits: your gain less the expected gain. Here $+0.67$ bits for the single
word, $-0.33$ for the pair. Added up over a game, this is the part of your bits that was luck.

\section{Risk}\label{sec:risk}

Some guesses gain about the same whatever the colours. Others usually gain a little and now and
then a lot: a gamble. Risk is the chance of gaining less than expected, minus the chance of gaining
more, as a percentage of the riskiest guess there was (never below 0).

In the example the chance of gaining less than $0.92$ is $2/3$ and of gaining more $1/3$, so the
raw risk is $+33$ points. A guess with equal groups has none. Risk belongs to the guess, not to
how it turned out, so it is the same however lucky you were.

\section{Your whole game}\label{sec:game}

The bits you gain guess by guess add up. When you find the word you have gained all the bits there
were to find. For a game: bits gained = bits expected (your skill) + luck. A game's picture is a
bar for each guess racing towards the bits needed.

\section{The last try}

On the last try nothing is left to learn: only a word that fits can win. Skill is the chance your
guess wins, against the likeliest word's. Luck is the outcome less that chance: if three words fit
and you guess one, a win is $+67$ and a miss is $-33$.

\section{Several words at once}

In a quartet one guess serves four boards. Each board is scored on its own, and the guess's skill
and luck are the averages over the boards that are still open, each counted by the bits it still
has to find.

\section{Playstyles}\label{sec:styles}

Your style is read from your last 100 games. Early on (64 or more words left) each guess counts
as the first that fits: a \textbf{gambler} (risk 75 or more), a \textbf{maximiser} (in the top
10\% of guesses by expected gain), an \textbf{eliminator} (four or more letters never tried).
At the end (16 or fewer words, or three greens known) each guess counts as a \textbf{closer} (it
fits every clue, so it could win) or a \textbf{prober} (it does not, ranks in the top 10\%, and
three or more words are left).

Your \textbf{habit} in a phase is what at least 40\% of its guesses were, and the most common if
several are. Your style is your early habit and your end habit:

\begin{center}
\begin{tabular}{lll}
\toprule
 & prober at the end & closer at the end\\
\midrule
maximiser early & analyst & maverick\\
eliminator early & detective & sweeper\\
gambler early & wildcard & gambler\\
\bottomrule
\end{tabular}
\end{center}

You are rated after ten games with ten guesses in each phase. These thresholds come from simulated
players, not from real ones, and may change.

\end{document}
